\chapter{卷积码方案与仿真分析}\label{chap:cc}

\section{高速电文卷积码设计目标}

本章面向长度固定为300 bit的高速二进制电文。与上一章讨论的低速电文相比，高速业务不仅要求降低整帧错误率，还要求编码器能够随数据到达连续工作，译码器能够在整帧结束前逐批给出稳定结果，并在发送效率、等待时间和存储资源之间作出可解释的取舍。若只追求最低帧错误率，可以增加冗余并等待完整码字；若只追求短发送时长，又会压缩纠错余量。因而本章不把某一个码率或某一组译码参数预设为唯一答案，而是先建立统一编码结构，再由正式仿真分别筛选可靠性、吞吐、时延、存储和综合平衡方案。

卷积码适合承担这一工程基线。编码输出由当前输入和有限长度历史共同决定，移位寄存器状态可以在数据连续到达时自然延续，不必把每个时隙重新初始化为独立码块。以码率$1/2$母码为基础，还可以通过打孔删除预定位置的编码比特，在不改变主体网格结构和Viterbi译码框架的条件下形成约$2/3$和约$3/4$方案。这种做法使可靠性与发送效率的比较具有共同基础，也便于接收端在硬判决和软判决之间切换。

接收侧采用Viterbi算法在状态网格中寻找累计路径度量最优的输入序列。整块译码能够利用完整终止边界，适合作为可靠性基准；有限回溯和真滑窗译码则用受控的历史长度、接收窗口和滑动步长换取更早输出与更低存储。对于按时隙到达的高速电文，连续编码与滑窗译码还可把“等待完整300 bit后统一处理”改为“边到达、边更新、分批提交”。

因此，本章围绕四个问题展开：第一，$K=7$、八进制生成多项式171/133如何适配300 bit电文并形成612、459和408 bit三种发送长度；第二，不同码率以及硬、软判决如何影响BER和FER；第三，回溯深度、软信息量化位宽及滑窗$W/S/D$参数如何改变可靠性与资源；第四，在最终码流不变的情况下，时隙组织为何保持BER/FER一致却产生不同输出时延。下文先从编码结构和实际码长开始。

评价这些问题时，FER用于回答一封300 bit电文能否完整交付，BER用于观察错误比特的总体密度，发送长度反映占用的调制符号数，首次输出和第95百分位决策时延则分别描述最早结果与大多数结果的等待。四类指标不能相互替代：较低BER并不必然保证整帧正确，较早出现第一个比特也不代表后续信息均已完成判决，而较高码率只有在给定信道余量下仍满足目标FER才具有业务价值。本章据此坚持在相同口径内比较，在目标改变时重新说明比较基准，避免用单一曲线或单一资源数值代替系统选型。

\section{卷积编码结构与码长设计}

\subsection{约束长度与生成多项式}

编码器采用约束长度$K=7$的二进制卷积母码。每个输出不仅由当前输入比特决定，还与前6 bit历史有关，故记忆长度为6，状态数为$2^6=64$。每输入1 bit，编码器沿状态网格前进一步并产生两个母码输出比特。两个输出分别由八进制生成多项式$G_1=171$和$G_2=133$规定；它们对应二进制多项式
\begin{equation}
  g_1(x)=x^6+x^5+x^4+x^3+1,\qquad
  g_2(x)=x^6+x^4+x^3+x+1 .
\end{equation}
这里的171和133是八进制表示，不是普通十进制系数。每个多项式标明当前输入和哪些记忆单元参加模2加法，两个异或网络由此产生一对编码比特。$K=7$使网格规模仍可由64状态Viterbi稳定处理，同时保留足够的路径记忆，是本项目300 bit高速电文的统一母码结构。

图\ref{fig:cc-encoder-placeholder}预留为编码器结构图。后续Visio绘制应完整表示当前输入、6级记忆单元、两组抽头和两个输出，不能只画抽象方框而省略多项式含义。

\begin{figure}[htbp]
  \centering
  \fbox{\begin{minipage}[c][0.21\textheight][c]{0.82\textwidth}
    \centering\Large 待插入：Visio绘制的卷积编码器移位寄存器与生成多项式结构图\\[0.8em]
    \normalsize 输入信息比特 $\rightarrow$ 当前输入与6级记忆单元 $\rightarrow$ 两个异或网络 $\rightarrow$ 两路编码输出\\[0.4em]
    标注：$K=7$，$G_1=171$（八进制），$G_2=133$（八进制）
  \end{minipage}}
  \caption{卷积编码器移位寄存器与生成多项式结构}
  \label{fig:cc-encoder-placeholder}
\end{figure}

\begin{table}[htbp]
  \centering
  \small
  \caption{卷积码基础编码参数}
  \label{tab:cc-basic-parameters}
  \begin{tabular}{p{0.32\textwidth}p{0.48\textwidth}}
    \toprule
    参数 & 冻结取值 \\
    \midrule
    原始信息长度 & 300 bit \\
    约束长度、记忆长度 & $K=7$；记忆长度6 bit \\
    状态数 & 64 \\
    生成多项式 & $G_1=171$、$G_2=133$（八进制） \\
    母码率 & $1/2$ \\
    尾比特 & 6 bit零尾 \\
    编码器输入长度 & 306个网格步 \\
    母码输出长度 & 612 bit \\
    \bottomrule
  \end{tabular}
\end{table}

\subsection{零尾终止与实际编码长度}

整块编码在300 bit原始电文后追加6 bit零尾，使编码器从任意数据结束状态逐步返回全零状态。编码器实际处理
\begin{equation}
  N_{\mathrm{step}}=300+6=306
\end{equation}
个网格步。母码每步输出2 bit，因此未打孔码长为
\begin{equation}
  N_{\mathrm{mother}}=306\times2=612~\mathrm{bit}.
\end{equation}
这6个尾步不属于原始业务信息，但为整块Viterbi提供已知终止状态。译码器可从零状态反向沿幸存路径回溯，避免把帧末状态当作未知量；恢复后只提交前300 bit原始电文。612 bit远低于1000 bit码长上限，也为后续提高码率保留了打孔空间。

零尾是整块组织的终止手段，不等于每个时隙都重新补6 bit。若把50、100或150 bit时隙各自独立终止，不仅会反复增加尾比特，还会切断跨时隙状态记忆，形成与本章正式方案不同的多个短卷积码块。连续时隙方案只在整封300 bit电文的最后一个时隙追加零尾，具体机制见第\ref{sec:cc-slot}节。

\subsection{打孔与实际码率}

母码输出按“每个网格步先输出$G_1$、再输出$G_2$”的顺序排列。打孔器周期性保留或删除这些位置：图样11保留全部母码比特；图样1101每4个母码位置保留3个；图样110110每6个位置保留4个。接收端解打孔时，被删除的位置以“无观测”标志送入路径度量计算，并不把它们错误补成判决比特0。因此，打孔改变的是可用观测数量和冗余，而不是强行引入固定比特。

实际帧级码率统一定义为
\begin{equation}
  R=\frac{K_{\mathrm{info}}}{N_{\mathrm{tx}}},
\end{equation}
其中$K_{\mathrm{info}}=300$，$N_{\mathrm{tx}}$是实际发送编码长度。表\ref{tab:cc-puncturing}中的“$1/2$、约$2/3$、约$3/4$”是母码或目标打孔方案名称；由于6个零尾不计入原始信息，实际帧级码率分别为0.4902、0.6536和0.7353，不能直接写成理想分数。

\begin{table}[htbp]
  \centering
  \small
  \caption{不同打孔方案的实际编码长度与码率}
  \label{tab:cc-puncturing}
  \begin{tabular}{p{0.19\textwidth}c c c p{0.23\textwidth}}
    \toprule
    方案 & 打孔图样 & 发送长度/bit & 实际码率 & 工程意义 \\
    \midrule
    $1/2$母码 & 11 & 612 & 0.4902 & 冗余最多，可靠性优先 \\
    约$2/3$打孔 & 1101 & 459 & 0.6536 & 可靠性与发送效率折中 \\
    约$3/4$打孔 & 110110 & 408 & 0.7353 & 发送长度最短，吞吐优先 \\
    \bottomrule
  \end{tabular}
\end{table}

\section{Viterbi译码与连续处理机制}

\subsection{网格状态、路径度量与回溯}

编码器的6 bit记忆内容构成网格状态。每到来1 bit输入，当前状态按寄存器移位规则转移到下一状态，并产生一对候选编码输出；这一过程称为一个网格步。每个状态有两条可能输入分支，连续分支连接成候选路径。接收端根据观测与候选输出的差异计算分支度量，再把它加到前一状态的累计路径度量上。到达同一新状态的多条路径中，仅保留累计度量较小者及其前驱信息，称为幸存路径。

Viterbi算法并不逐个比特作孤立判决，而是比较完整状态约束下的路径。噪声可能使某一步的局部观测倾向错误分支，但后续观测继续累积后，不符合卷积约束的路径往往付出更大度量而被淘汰。经过足够深度后，从当前优胜状态向前回溯，早期路径通常已经合并，此时才能把对应信息比特作为稳定结果提交。图\ref{fig:cc-trellis-placeholder}将用少量状态示意分支竞争和幸存路径；实际$K=7$译码器始终包含64状态，示意图中的状态数量不代表真实规模。

路径合并也是理解回溯和滑窗参数的关键。若两条候选路径在较早时刻已经进入同一状态，则该状态之前的共同部分不会再被后续输入分开；反之，若多条竞争路径在有限历史内仍保持接近的累计度量，立即输出就可能选错。回溯深度给竞争路径留下继续分化的观察时间，窗口长度提供可用于比较的接收范围，滑动步长决定稳定结果以多大批次被提交。三者分别作用于路径稳定、观察范围和输出节奏，不能把其中一个参数的增大简单视为对另外两个参数的替代。

\begin{figure}[htbp]
  \centering
  \fbox{\begin{minipage}[c][0.21\textheight][c]{0.82\textwidth}
    \centering\Large 待插入：Visio绘制的网格状态转移与Viterbi幸存路径示意图\\[0.8em]
    \normalsize 少量状态仅用于说明分支、累计路径度量、路径竞争与回溯；实际译码器共有64状态
  \end{minipage}}
  \caption{网格状态转移与Viterbi幸存路径示意}
  \label{fig:cc-trellis-placeholder}
\end{figure}

\subsection{硬判决与软判决}

硬判决先把BPSK接收样值量化为0或1，再统计候选编码输出与接收判决之间不一致的位置数，形成汉明型分支度量。该接口简单，每个观测只需1 bit，但量化后无法区分“接收点刚刚越过判决门限”和“接收点远离门限”这两种可靠性完全不同的情况。

浮点软判决保留BPSK接收样值。对每个候选编码输出，先映射到理想符号$+1$或$-1$，再累计接收样值与候选符号之间的平方欧氏距离。当前实现不是直接输入对数似然比（LLR），而是用连续接收样值中的幅度距离表达可靠性。对于打孔删除的位置，观测标志为0，该位置不参与硬判决或软判决的分支度量。软判决因此没有改变卷积码和状态网格，只是用更完整的信道信息提高了路径区分能力。

\subsection{整块、有限回溯与滑窗译码}

整块译码接收完整306步网格，保存全帧幸存路径，从已知零终止状态回溯并一次输出300 bit。它最大限度利用完整帧边界，是比较码率和判决方式时的可靠性基准，但需要保存约306步的幸存信息，且最早结果必须等到帧末。

有限回溯仍持续更新64个状态度量，但只保留最近$D$步路径历史；当路径足够稳定时提交较早比特。增大$D$通常有助于降低过早判决风险，却同步增加存储和等待，并可能带来重复回溯操作。真滑窗在此基础上引入接收窗口长度$W$和滑动步长$S$：窗口缓存连续观测，更新路径度量，向前回溯$D$步，每个非最终窗口提交约$S$个稳定信息比特，然后前移$S$继续处理。参数必须满足
\begin{equation}
  W>D,\qquad S\leq W-D.
\end{equation}
最终窗口利用零终止状态完成剩余输出，并检查300 bit是否无遗漏、无重复。这一过程是真正的增量译码，不是先做整帧译码再人为截断结果。图\ref{fig:cc-block-sliding-placeholder}预留两条流程，后续应明确整块一次输出与滑窗分批提交的差别。

\begin{figure}[htbp]
  \centering
  \fbox{\begin{minipage}[c][0.23\textheight][c]{0.84\textwidth}
    \centering\Large 待插入：Visio绘制的整块与滑窗Viterbi译码流程图\\[0.8em]
    \normalsize 整块：整帧接收 $\rightarrow$ 完整路径度量 $\rightarrow$ 完整幸存路径 $\rightarrow$ 终止状态回溯 $\rightarrow$ 一次输出300 bit\\[0.5em]
    滑窗：连续接收 $\rightarrow$ 缓存$W$ $\rightarrow$ 度量更新 $\rightarrow$ 回溯$D$ $\rightarrow$ 提交$S$个稳定比特 $\rightarrow$ 前移$S$
  \end{minipage}}
  \caption{整块与滑窗Viterbi译码流程}
  \label{fig:cc-block-sliding-placeholder}
\end{figure}

\section{码率与判决方式性能分析}

\subsection{仿真条件}

基础性能比较采用300 bit独立等概率二进制电文、BPSK调制和实值AWGN信道。主横轴为$E_s/N_0$，单位dB；幅度为$\pm1$时，噪声方差为
\begin{equation}
  \sigma^2=\frac{1}{2\times10^{(E_s/N_0)/10}}.
\end{equation}
若需要换算为$E_b/N_0$，应对每个方案分别使用$E_b/N_0=E_s/N_0-10\log_{10}R$，其中$R$取表\ref{tab:cc-puncturing}的实际帧级码率，不能对三码率统一平移。

\begin{table}[htbp]
  \centering
  \small
  \caption{基础性能仿真参数}
  \label{tab:cc-simulation-settings}
  \begin{tabular}{p{0.28\textwidth}p{0.55\textwidth}}
    \toprule
    项目 & 设置 \\
    \midrule
    原始电文与编码 & 300 bit；$K=7$；171/133（八进制）；6 bit零尾 \\
    调制与信道 & BPSK；实值AWGN \\
    主横轴 & $E_s/N_0$ \\
    基础粗网格 & $-5\sim10$ dB，步长0.5 dB \\
    粗网格停止条件 & 每点至少1000帧；累计200个错误帧后可停止；最多50000帧 \\
    密集区用途 & 每点至少5000帧，用于瀑布区及FER=0.1工作点插值 \\
    码率方案 & $1/2$母码、约$2/3$打孔、约$3/4$打孔 \\
    判决方式 & 硬判决、浮点软判决 \\
    主要指标 & BER、FER；资源与时延另行统计 \\
    \bottomrule
  \end{tabular}
\end{table}

粗网格用于覆盖完整工作区间，密集区数据用于提高瀑布段工作点插值的分辨率，两者停止条件不能混写。FER=0.1所需信噪比在相邻非零统计点间对$\log_{10}(\mathrm{FER})$线性插值。若高信噪比点在有限帧数内未观察到错误，只能表述为“零误帧观测”；在对数图中省略该点，不把它解释为理论FER等于零。

\subsection{不同码率可靠性分析}

图\ref{fig:cc-rate-soft-fer}在整块浮点软判决条件下比较三码率。提高码率会减少发送冗余和实际码长：约$3/4$方案只发送408 bit，约$2/3$方案发送459 bit，$1/2$母码发送612 bit。相应地，在相同$E_s/N_0$下，打孔方案可供路径度量使用的观测更少，瀑布区依次向右移动。

\reportfigure[0.88\textwidth]{figures/cc/cc_rate_soft_fer_cn.png}
  {不同码率浮点软判决的帧错误率}{fig:cc-rate-soft-fer}
  {横轴为$E_s/N_0$；有限样本零误帧点未在对数坐标中绘制。}

\begin{table}[htbp]
  \centering
  \small
  \caption{不同码率的FER=0.1工作点}
  \label{tab:cc-rate-fer01}
  \begin{tabular}{p{0.25\textwidth}c c p{0.30\textwidth}}
    \toprule
    方案 & 发送长度/bit & 所需$E_s/N_0$/dB & 工程解释 \\
    \midrule
    $1/2$母码 & 612 & $-0.709$ & 冗余最多，可靠性余量最大 \\
    约$2/3$打孔 & 459 & $0.963$ & 减少153 bit，形成折中 \\
    约$3/4$打孔 & 408 & $1.885$ & 发送最短，以可靠性换取效率 \\
    \bottomrule
  \end{tabular}
\end{table}

表\ref{tab:cc-rate-fer01}量化了这一取舍。相对$1/2$母码，约$2/3$和约$3/4$方案达到FER=0.1分别需要更高的$E_s/N_0$，但不能据此简单称约$3/4$方案“最差”：其代价对应204 bit的码长缩短和更高实际码率，适用于链路余量充足、发送占用或吞吐优先的场合。码率选择的实质是在可用信噪比下决定把资源留给冗余还是业务传输。

这里的公平性边界必须明确：比较固定的是每个实际发送BPSK符号的$E_s/N_0$。$1/2$母码、约$2/3$和约$3/4$方案每帧分别发送612、459和408个符号，故低码率方案不仅冗余更多，整帧占用的发送符号也更多。本节回答的是固定单符号能量条件下的“发送占用---可靠性”权衡，而不是固定每信息比特能量或固定整帧总能量条件下的纯编码增益。

\subsection{硬判决与软判决性能分析}

图\ref{fig:cc-hard-soft-fer}显示，对三个码率，浮点软判决瀑布区均明显早于硬判决进入下降段。两种译码器采用相同卷积码、打孔图样和状态转移，差别只在分支度量：硬判决丢弃门限两侧的距离，软判决保留接收点相对$\pm1$理想符号的连续幅度信息，因此能更早排除不可信路径。

\reportfigure[0.98\textwidth]{figures/cc/cc_hard_soft_fer_cn.png}
  {硬判决与浮点软判决的帧错误率比较}{fig:cc-hard-soft-fer}
  {三幅子图分别对应$1/2$母码、约$2/3$和约$3/4$打孔方案。}

\begin{table}[htbp]
  \centering
  \small
  \caption{软判决相对硬判决的FER=0.1性能收益}
  \label{tab:cc-soft-gain}
  \begin{tabular}{p{0.28\textwidth}c p{0.42\textwidth}}
    \toprule
    方案 & $E_s/N_0$收益/dB & 解释 \\
    \midrule
    $1/2$母码 & 2.085 & 利用幅度可靠性后所需信噪比显著下降 \\
    约$2/3$打孔 & 1.927 & 打孔观测减少，但软信息优势仍保持 \\
    约$3/4$打孔 & 1.857 & 在最高码率下仍获得明显收益 \\
    \bottomrule
  \end{tabular}
\end{table}

在FER=0.1处，软判决相对硬判决的$E_s/N_0$收益为1.857--2.085 dB。它不需要增加发送冗余，而是以更精细的接收输入、路径度量运算和存储精度换取可靠性。因此，后续回溯、量化、滑窗和时隙研究均以软判决作为主要性能路线，同时保留硬判决作为低接口复杂度基准。

\section{回溯深度与软信息量化}

\subsection{回溯深度对可靠性的影响}

完整回溯保存全帧306步幸存路径并利用零尾终止状态；有限回溯正式比较$D=35,49,70,84,98,112$。图\ref{fig:cc-traceback-fer}给出三码率的全部正式FER点。随着$D$增加，过早提交尚未合并路径的风险下降，各有限深度曲线逐步逼近完整回溯，而不是只有$D=112$一个孤立结果。

\reportfigure[0.99\textwidth]{figures/cc/round06_d/cc_traceback_fer_cn.png}
  {不同回溯深度下的卷积码帧错误率}{fig:cc-traceback-fer}
  {三幅子图分别对应三码率；零误帧观测点未在对数坐标中伪造。}

\subsection{回溯深度对时延、存储和软件处理的影响}

图\ref{fig:cc-traceback-cpu}是当前处理器、编译器和串行实现上的软件处理时间，并非硬件固有时延。符号级首次输出时延则由译码结构决定，是更稳定的算法指标。有限回溯还会重复执行局部回溯，因此深度缩短并不保证本机软件时间同步下降。

\reportfigure[0.82\textwidth]{figures/cc/round06_d/cc_traceback_cpu_cn.png}
  {不同回溯深度下的软件译码时间}{fig:cc-traceback-cpu}
  {柱高为三码率与三个目标工作点的平均值，仅适用于当前软件平台。}

\reportfigure[0.96\textwidth]{figures/cc/round06_d/cc_traceback_tradeoff_cn.png}
  {有限回溯深度的可靠性与资源权衡}{fig:cc-traceback-tradeoff}
  {左图将全工作点最坏值与5\%判据直接比较；中位数仅作辅助观察。}

\begin{table}[htbp]
  \centering
  \scriptsize
  \caption{不同回溯深度的可靠性、资源与时延比较}
  \label{tab:cc-traceback-results}
  \begin{tabular}{c c c c c c p{0.22\textwidth}}
    \toprule
    $D$ & 最坏FER增幅 & 5\%判据 & 存储/B & 首次输出/符号 & 软件时间/$\mu$s & 评价 \\
    \midrule
    35  & 445\% & 不通过 & 7744  & 34  & 315.919 & 可靠性损失很大 \\
    49  & 128\% & 不通过 & 10432 & 48  & 319.674 & 与完整回溯差距大 \\
    70  & 29.1\% & 不通过 & 14464 & 69  & 329.360 & 可靠性继续改善 \\
    84  & 12.5\% & 不通过 & 17152 & 83  & 334.853 & 资源较低，未过门限 \\
    98  & 6.76\% & 不通过 & 19840 & 97  & 341.030 & 略高于5\%门限 \\
    112 & 3.38\% & 通过   & 22528 & 111 & 首个达标有限深度 \\
    完整 & 基准 & 基准 & 59776 & 305 & 288.388 & 全帧可靠性基准 \\
    \bottomrule
  \end{tabular}
\end{table}

$D=112$是当前测试集合中第一个通过全工作点5\%最坏相对FER增幅门限的有限深度；$D=98$仍略高于门限。相对完整回溯，$D=112$将译码存储降低约62.3\%，首次输出提前194个符号，故通用选择保持$D=112$。

\subsection{不同量化位宽的完整FER性能}

正式试验覆盖浮点以及3--8 bit六种量化输入。图\ref{fig:cc-quantization-fer}直接展示全部逐点FER曲线，未用插值生成曲线点；高信噪比的零误帧观测在对数轴上省略。

\reportfigure[0.99\textwidth]{figures/cc/round06_d/cc_quantization_fer_cn.png}
  {不同软信息量化位宽下的卷积码帧错误率}{fig:cc-quantization-fer}
  {位宽图例中的3--8 bit后文简称Q3--Q8。}

\begin{table}[htbp]
  \centering
  \small
  \caption{软信息量化实现参数}
  \label{tab:cc-quantization-rule}
  \begin{tabular}{p{0.18\textwidth}p{0.72\textwidth}}
    \toprule
    项目 & 正式实现 \\
    \midrule
    位宽含义 & Q$b$为$b$ bit有符号输入码，$q_{\max}=2^{b-1}-1$，使用对称整数区间$[-q_{\max},q_{\max}]$ \\
    输入与缩放 & 对解打孔后的实值接收样值量化；裁剪幅度Q3--Q6为1.5，Q7--Q8为2.0；步长$\Delta=A/q_{\max}$ \\
    舍入与饱和 & 先按$\operatorname{lround}(y/\Delta)$舍入到最近整数，再限幅到$[-q_{\max},q_{\max}]$ \\
    路径度量 & 量化输入采用平方欧氏分支度量，累计量为32 bit有符号整数；每个网格步减去最小度量进行归一化 \\
    溢出策略 & 候选值以64 bit临时量计算，超出32 bit或$10^9$上限时饱和；正式结果的两类计数均为0 \\
    研究边界 & 不仅输入被量化，量化译码支路的路径度量也为整数；浮点基准另用浮点软译码器 \\
    \bottomrule
  \end{tabular}
\end{table}

\subsection{量化性能损失与位宽选择}

图\ref{fig:cc-quantization-loss}在FER=0.1处比较量化相对浮点的信噪比差值，逻辑顺序是“完整曲线---工作点差值---工程选择”。

\reportfigure[0.86\textwidth]{figures/cc/round06_d/cc_quantization_loss_cn.png}
  {软信息量化位宽的FER=0.1性能损失}{fig:cc-quantization-loss}
  {负的微小差值来自有限样本及工作点插值波动，不表示量化理论性能优于浮点。}

\begin{table}[htbp]
  \centering
  \small
  \caption{代表性软信息精度的FER=0.1性能损失与用途}
  \label{tab:cc-quantization-results}
  \begin{tabular}{p{0.13\textwidth}p{0.31\textwidth}p{0.21\textwidth}p{0.25\textwidth}}
    \toprule
    精度 & 三码率相对浮点损失/dB & 表示代价 & 推荐用途 \\
    \midrule
    浮点 & 0（基准） & 输入与度量精度最高 & 系统级软件基准 \\
    Q4 & 0.103 / 0.081 / 0.061 & 4 bit输入与整数度量 & 存储受限候选 \\
    Q6 & 0.071 / 0.013 / 0.000 & 6 bit输入与整数度量 & 中等位宽候选 \\
    Q8 & 0.006 / 0.008 / $-0.001$ & 8 bit输入与整数度量 & 量化实现推荐 \\
    \bottomrule
  \end{tabular}
\end{table}

Q8在三码率上的损失分别为0.006214、0.008130和$-0.000741$ dB，已接近统计与插值分辨率，因此量化实现推荐仍为Q8；系统级五类方案仍以浮点软判决为正式比较口径。

\section{滑窗译码参数分析}

\subsection{窗口长度影响}

窗口长度试验取$W\in\{96,128,160,192\}$，固定$S=16,D=70$。图\ref{fig:cc-sliding-w}表明，在当前控制条件下各曲线已高度接近：$W$达到足够观察范围后，继续增加窗口的可靠性收益趋于饱和，却仍增加缓存和首窗口等待。

\reportfigure[0.99\textwidth]{figures/cc/round06_d/cc_sliding_w_fer_cn.png}
  {不同窗口长度对滑窗译码FER的影响}{fig:cc-sliding-w}
  {仅改变$W$，固定$S=16,D=70$。}

\subsection{滑动步长影响}

步长试验取$S\in\{8,16,25,50\}$，固定$W=160,D=70$。$S$主要控制输出批量和触发频率：较小$S$触发更频繁、重复计算更多；较大$S$扩大单批提交量并改变等待分布。它不是纠错深度，不能把“$S$越大”写成“纠错能力越强”。

\reportfigure[0.99\textwidth]{figures/cc/round06_d/cc_sliding_s_fer_cn.png}
  {不同滑动步长对滑窗译码FER的影响}{fig:cc-sliding-s}
  {仅改变$S$，固定$W=160,D=70$。}

\subsection{回溯深度影响}

滑窗回溯试验取$D\in\{35,49,70,84,98,112\}$，固定$W=160,S=16$。小$D$增加过早提交风险，增大$D$后可靠性逐渐逼近饱和，同时提高等待、幸存路径存储和回溯工作量。约$3/4$最终候选的$D=126$没有进入这组完整控制扫描，只参加最终候选比较。

\reportfigure[0.99\textwidth]{figures/cc/round06_d/cc_sliding_d_fer_cn.png}
  {不同回溯深度对滑窗译码FER的影响}{fig:cc-sliding-d}
  {仅改变$D$，固定$W=160,S=16$；不含最终候选专用的$D=126$。}

\begin{table}[htbp]
  \centering
  \scriptsize
  \caption{滑窗控制变量试验摘要}
  \label{tab:cc-sliding-controls}
  \begin{tabular}{p{0.08\textwidth}p{0.17\textwidth}p{0.13\textwidth}p{0.20\textwidth}p{0.20\textwidth}}
    \toprule
    变量 & 正式集合 & 固定条件 & 结构统计 & 工程解释 \\
    \midrule
    $W$ & 96/128/160/192 & $S=16,D=70$ & 首次输出：170/127/113符号 & 观察充分后FER收益趋于饱和 \\
    $S$ & 8/16/25/50 & $W=160,D=70$ & $S=16$时P95：260/195/173符号 & 控制输出节奏与触发频率 \\
    $D$ & 35，49，70，84，98，112 & $W=160,S=16$ & 逐点曲线覆盖六个深度 & 增深降低过早提交风险并增加等待、存储 \\
    \bottomrule
  \end{tabular}
\end{table}

\subsection{滑窗参数综合选择}

筛选遵循先满足可靠性、再比较等待、存储和处理量的顺序，不对$W/S/D$机械取最大或最小。

\begin{table}[htbp]
  \centering
  \scriptsize
  \caption{不同码率和目标下的滑窗参数推荐}
  \label{tab:cc-sliding-recommendations}
  \begin{tabular}{p{0.17\textwidth}p{0.19\textwidth}c c c p{0.31\textwidth}}
    \toprule
    码率方案 & 目标 & $W$ & $S$ & $D$ & 选择含义 \\
    \midrule
    $1/2$母码 & 综合/性能 & 128 & 25 & 70 & 冗余充足，较浅回溯即可稳定 \\
    $1/2$母码 & 时延/存储 & 96 & 16 & 70 & 缩小缓存并保持可接受可靠性 \\
    约$2/3$打孔 & 综合/性能 & 128 & 25 & 98 & 以更深回溯补偿观测减少 \\
    约$2/3$打孔 & 时延/存储 & 128 & 25 & 84 & 在等待与可靠性间折中 \\
    约$3/4$打孔 & 综合/性能 & 160 & 25 & 126 & 最少冗余需要更长观察与回溯 \\
    约$3/4$打孔 & 时延/存储 & 128 & 25 & 98 & 接受性能代价以减少资源 \\
    \bottomrule
  \end{tabular}
\end{table}

不存在跨三码率统一最优的$W/S/D$。打孔越多，竞争路径通常需要更长观察或回溯；目标函数不同，也会在相同码率下给出不同配置。

\section{时隙组织与连续译码性能}\label{sec:cc-slot}

\subsection{连续编码与时隙组织方式}

正式比较包括整块300 bit，以及$50\times6$、$100\times3$、$150\times2$三种连续时隙组织。连续方案跨时隙保持编码器状态和打孔相位，只有最后一个时隙追加6 bit零尾；三种连续组织编码的是同一封300 bit电文，而不是多个独立短码块。

\begin{figure}[htbp]
  \centering
  \fbox{\begin{minipage}[c][0.20\textheight][c]{0.86\textwidth}
    \centering\Large 待插入：Visio绘制的高速电文按时隙连续编码与滑窗译码流程图\\[0.8em]
    \small 300 bit高速电文 $\rightarrow$ 按时隙拆分 $\rightarrow$ 连续卷积编码 $\rightarrow$ BPSK $\rightarrow$ AWGN\\[0.4em]
    $\rightarrow$ 按时隙到达 $\rightarrow$ 接收缓存 $\rightarrow$ 滑窗Viterbi $\rightarrow$ 逐批恢复信息比特\\[0.5em]
    标注：编码器状态与打孔相位跨时隙保持，仅最后一个时隙执行零尾终止
  \end{minipage}}
  \caption{高速电文按时隙连续编码与滑窗译码流程}
  \label{fig:cc-slot-flow-placeholder}
\end{figure}

\subsection{硬判决下的时隙可靠性}

图\ref{fig:cc-slot-hard-fer}比较整块与三种连续组织的硬判决FER。对应BER组合图也已按同一正式CSV生成，并为控制正文篇幅移至附录\ref{app:figures}；300 bit电文的主要交付指标采用FER。

\reportfigure[0.99\textwidth]{figures/cc/round06_d/cc_slot_hard_fer_cn.png}
  {不同电文组织的硬判决帧错误率}{fig:cc-slot-hard-fer}
  {三幅子图分别对应三码率；每图比较整块300和三种连续组织。}

\subsection{浮点软判决下的时隙可靠性}

图\ref{fig:cc-slot-soft-fer}给出浮点软判决FER。相同码率、判决方式和$E_s/N_0$下，$50\times6$、$100\times3$和$150\times2$三种连续组织的BER与FER逐点完全重合。原因是最终发送码流、编码器状态、打孔相位、仅末时隙终止、接收序列、噪声、滑窗规则均一致；时隙组织改变的是接收数据何时到达，而不是最终卷积码字。不能据此声称$50\times6$纠错更强。

\reportfigure[0.99\textwidth]{figures/cc/round06_d/cc_slot_soft_fer_cn.png}
  {不同电文组织的浮点软判决帧错误率}{fig:cc-slot-soft-fer}
  {三种连续组织曲线逐点重合；软判决BER组合图见附录\ref{app:figures}。}

\subsection{不同时隙组织的输出时延}

时延以接收符号计数。首次输出反映最早结果，第95百分位反映绝大多数信息比特的等待，整帧最后决策符号反映全部交付时刻。

\reportfigure[0.98\textwidth]{figures/cc/round06_d/cc_slot_latency_cn.png}
  {不同时隙组织的浮点软判决输出时延}{fig:cc-slot-latency}
  {首次输出与第95百分位描述不同的等待特征。}

\begin{table}[htbp]
  \centering
  \scriptsize
  \caption{不同时隙组织的浮点软判决输出时延}
  \label{tab:cc-slot-latency}
  \begin{tabular}{c p{0.13\textwidth}c c c c c}
    \toprule
    码率 & 组织方式 & 首次 & 平均 & P95 & 最大 & 整帧最后符号 \\
    \midrule
    $1/2$ & 整块300 & 610 & 311.0 & 580 & 610 & 611 \\
    $1/2$ & $50\times6$ & 298 & 231.0 & 342 & 360 & 611 \\
    $1/2$ & $100\times3$ & 398 & 258.0 & 430 & 460 & 611 \\
    $1/2$ & $150\times2$ & 298 & 285.0 & 530 & 560 & 611 \\
    \midrule
    约$2/3$ & 整块300 & 457 & 233.0 & 435 & 457 & 458 \\
    约$2/3$ & $50\times6$ & 223 & 173.0 & 256 & 270 & 458 \\
    约$2/3$ & $100\times3$ & 298 & 193.25 & 322 & 345 & 458 \\
    约$2/3$ & $150\times2$ & 223 & 213.5 & 397 & 420 & 458 \\
    \midrule
    约$3/4$ & 整块300 & 406 & 207.0 & 386 & 406 & 407 \\
    约$3/4$ & $50\times6$ & 265 & 171.17 & 262 & 273 & 407 \\
    约$3/4$ & $100\times3$ & 265 & 183.5 & 320 & 340 & 407 \\
    约$3/4$ & $150\times2$ & 406 & 207.0 & 386 & 406 & 407 \\
    \bottomrule
  \end{tabular}
\end{table}

约$2/3$软判决的$50\times6$和$150\times2$均在223个符号后首次输出，但P95分别为256和397个符号。$50\times6$的低时延结论来自整体等待分布，而不是可靠性差异。

\subsection{不同时隙组织的有效吞吐率}

归一化有效吞吐率同时受实际码率与FER影响。高码率并不自动产生更高成功业务交付；只有可靠性仍可接受时，较短发送长度才会转化为有效交付优势。

\reportfigure[0.99\textwidth]{figures/cc/round06_d/cc_slot_soft_goodput_cn.png}
  {不同电文组织的浮点软判决归一化有效吞吐率}{fig:cc-slot-goodput}
  {曲线来自正式逐点数据；该指标是有效吞吐率，不称为编码增益。}

硬判决有效吞吐率图已有正式数据支撑，但与硬判决FER图传达的信息重复，故不进入正文；正文保留浮点软判决主路线。

\subsection{时隙组织选择}

在三种连续组织可靠性完全相同的前提下，选择由到达节奏和时延分布决定。$50\times6$更频繁地向接收缓存提供新符号，在约$2/3$方案中保持223符号的最早首次输出，并把P95压缩到256符号，因此继续作为低时延推荐。该结论依赖跨时隙状态与打孔相位连续；若各时隙独立清零或重新起相位，必须重新评估。

\section{卷积码方案阶段性结论}

表\ref{tab:cc-main-conclusions}把本章各控制问题压缩为可执行选择。它们共同构成一条完整工程链：300 bit电文采用$K=7$、171/133母码和6 bit零尾；通过打孔获得612、459和408 bit三种发送长度；浮点软判决建立性能基准；有限回溯、量化和滑窗再按资源目标降维；连续时隙在保持可靠性的同时调整交付时延。

\begin{table}[htbp]
  \centering
  \small
  \caption{卷积码主要参数选择结论}
  \label{tab:cc-main-conclusions}
  \begin{tabular}{p{0.22\textwidth}p{0.25\textwidth}p{0.43\textwidth}}
    \toprule
    设计问题 & 最终结论 & 工程原因 \\
    \midrule
    可靠性码率 & $1/2$母码 & 冗余最多，FER门限最低 \\
    吞吐码率 & 约$3/4$打孔 & 408 bit发送长度，实际码率最高 \\
    判决方式 & 浮点软判决 & FER=0.1处获得约1.86--2.08 dB收益 \\
    通用有限回溯 & $D=112$ & 唯一通过5\%最坏相对FER增幅判据，存储降约62.3\% \\
    量化实现 & Q8 & 三码率相对浮点损失不超过约0.008 dB量级 \\
    滑窗配置 & 按码率和目标分别设置 & 冗余度改变路径收敛，时延与存储目标也不同 \\
    连续时隙 & 低时延优先$50\times6$ & 可靠性与其他连续组织相同，第95百分位等待更低 \\
    \bottomrule
  \end{tabular}
\end{table}

\begin{table}[htbp]
  \centering
  \scriptsize
  \caption{卷积码五类系统级推荐}
  \label{tab:cc-system-recommendations}
  \begin{tabular}{p{0.15\textwidth}p{0.38\textwidth}p{0.20\textwidth}p{0.20\textwidth}}
    \toprule
    目标 & 推荐方案 & 比较基准 & 选择理由 \\
    \midrule
    可靠性优先 & $1/2$母码整块浮点软判决 & 固定$E_s/N_0=2$ dB & 共同信噪比下FER最低 \\
    吞吐优先 & 约$3/4$打孔滑窗，$W=160,S=25,D=126$ & 固定$E_s/N_0=2$ dB & 实际码率与有效吞吐率最高 \\
    时延优先 & 约$2/3$打孔滑窗，$W=128,S=25,D=84$ & 固定FER=0.1 & 覆盖FER门限且首次输出最早 \\
    存储优先 & $1/2$母码滑窗，$W=96,S=16,D=70$ & 固定FER=0.1 & 覆盖FER门限且译码存储最低 \\
    综合平衡 & $1/2$母码滑窗，$W=96,S=16,D=70$ & 固定FER=0.1 & 综合四项指标最优 \\
    \bottomrule
  \end{tabular}
\end{table}

五类推荐不是一个“统一最优方案”。可靠性与吞吐方案在固定$E_s/N_0$下比较，时延、存储和综合方案在固定FER=0.1下比较；基准不同，不能声称其中某一项全面优于其余方案。类似地，Q8量化建议回答的是定点输入位宽问题，不改变表\ref{tab:cc-system-recommendations}以浮点软判决建立的系统级推荐口径。

工程部署时还应把这些推荐视为可复核的起点，而不是脱离链路条件的永久常数。若业务信噪比裕量降低，应优先检查高码率方案是否仍覆盖目标FER；若接收缓存成为主要约束，可在同一码率内选择较小窗口配置；若处理器或硬件采用并行回溯，软件测得的时间排序也需要重新验证。尽管绝对资源会随平台变化，本章冻结的码长、状态数、路径度量、FER工作点和符号时延仍提供了统一接口，使后续比较能够明确区分“编码方案差异”和“实现平台差异”。

本章已经完成卷积码内部设计和筛选。后续横向比较将使用冻结的$1/2$母码612 bit与约$2/3$打孔459 bit整块浮点软判决候选，并采用完整306步终止回溯；复杂度分析和突发错误交织专题也分别沿用已冻结配置。此处只给出输入边界，不提前展开多信道结果、译码算法横向复杂度或交织收益。由此，第5章回答了“为什么这样设计、正式结果如何、差异为何出现以及不同工程目标如何选择”，并为后续系统级比较提供一致的高速卷积码基线。
